\chapter{Conclusão}
\label{cap:conclusao}

Este trabalho partiu do reconhecimento de um problema central na aplicação de inteligência artificial ao setor público: a tensão entre a performance de modelos de caixa-preta e a necessidade de motivação, auditabilidade e transparência exigidas pela atuação administrativa. No contexto do Rating da Dívida Ativa da SEFAZ-DF, esse problema assume especial relevância, uma vez que a classificação de recuperabilidade dos créditos influencia diretamente a priorização de esforços de cobrança e, por consequência, a própria racionalidade da gestão fiscal.

Com base nessa problemática, a pesquisa analisou a aplicação das técnicas SHAP e LIME sobre o modelo de produção utilizado no Rating da Dívida Ativa, preservando a lógica operacional do sistema por meio da reprodução do pré-processamento, do \textit{ensemble} de modelos XGBoost e das regras de pós-processamento. Em razão das restrições de sigilo e da remoção dos dados reais da versão final do projeto, a camada explicativa foi construída exclusivamente sobre bases sintéticas, sem prejuízo da aderência ao comportamento decisório do modelo original.

Os resultados mostraram que a proposta é tecnicamente viável e analiticamente útil. Em primeiro lugar, observou-se que o SHAP forneceu explicações globais e locais consistentes para o modelo analisado, permitindo identificar, com coerência, as variáveis mais relevantes para a classificação das CDAs. Em segundo lugar, verificou-se que o LIME também produziu explicações locais informativas, mas com maior sensibilidade ao contexto e maior variabilidade entre execuções, sobretudo nas posições inferiores do \textit{ranking} explicativo. Em terceiro lugar, a comparação entre os dois métodos revelou convergência parcial em torno de um núcleo recorrente de fatores, especialmente variáveis associadas ao histórico de regularização, à exposição do contribuinte no núcleo e ao tempo sem movimentação da dívida. Esse comportamento reforça a plausibilidade substantiva do modelo e, ao mesmo tempo, evidencia que a explicabilidade não deve ser tratada como propriedade automática de qualquer técnica interpretativa.

À luz desses achados, conclui-se que o SHAP se mostrou o método mais adequado para sustentar uma camada formal de explicabilidade no contexto da auditoria fiscal. Sua capacidade de oferecer interpretação global e local, aliada à consistência observada no \textit{ensemble} de produção, faz dele a alternativa mais robusta para apoiar processos de governança algorítmica e fundamentação administrativa. O LIME, por sua vez, mostrou-se valioso como ferramenta complementar de inspeção local e de comunicação com públicos não técnicos, mas não apresentou o mesmo nível de estabilidade para servir isoladamente como fundamento principal de motivação técnico-administrativa.

Do ponto de vista das contribuições do trabalho, destacam-se cinco resultados principais. O primeiro foi a demonstração de que é possível explicar o modelo real de produção, em vez de um modelo substituto simplificado, preservando a aderência à lógica decisória efetivamente utilizada pela SEFAZ-DF. O segundo foi a consolidação de uma estratégia metodológica baseada em dados sintéticos, capaz de compatibilizar restrições de sigilo com a necessidade de interpretação do modelo. O terceiro foi a comparação experimental entre SHAP e LIME em diferentes \textit{setups} e classes de \textit{rating}, oferecendo uma leitura mais criteriosa sobre consistência, concordância e estabilidade das explicações. O quarto foi a organização dos resultados em artefatos visuais e tabulares com potencial de apoiar rotinas de auditoria fiscal mais transparentes. O quinto foi a demonstração experimental, com um modelo substituto, de que um histórico que sub-registra os pagamentos de um subgrupo de contribuintes gera viés contra esse grupo sem alterar de forma perceptível a performance média do modelo, e de que o SHAP torna esse viés visível ao auditor.

Como limitações, o estudo deve ser interpretado à luz de quatro restrições. A primeira é que a versão final da pesquisa não utiliza dados reais, o que impede a reaplicação integral de certas etapas exploratórias originalmente concebidas para comparação direta entre bases reais e sintéticas. A segunda decorre do próprio sintetizador: a cópula gaussiana preserva as distribuições e as correlações entre variáveis, mas não reproduz com fidelidade a relação entre as características da CDA e o pagamento observado, de modo que as explicações descrevem como o modelo decide diante de casos verossímeis, e não o comportamento real dos contribuintes; além disso, a proporção de CDAs parceladas ou suspensas nas Taxas não foi bem preservada, o que afeta a regra de \textit{rating} A forçado nesse imposto. A terceira é a assimetria de maturidade entre os \textit{setups}, uma vez que o \textit{ICMS} concentrou o conjunto mais completo de artefatos visuais e analíticos. A quarta é que, embora a camada explicativa tenha sido aplicada ao modelo de produção, sua operacionalização final em ambiente institucional dependeria de integração formal com os fluxos de trabalho e critérios de governança da Administração Tributária.

Ainda assim, tais limitações não comprometem a contribuição principal da pesquisa. Ao demonstrar que o modelo real de Rating da Dívida Ativa pode ser interpretado de forma consistente por meio de SHAP, com apoio complementar do LIME e sem exposição de dados reais, o trabalho oferece uma resposta concreta ao dilema entre a performance do modelo e a transparência. Nesse sentido, a principal conclusão desta monografia é que a explicabilidade do Rating não apenas é tecnicamente possível, mas constitui um passo necessário para ampliar a auditabilidade, a legitimidade e a segurança jurídica da aplicação de inteligência artificial na gestão fiscal.
